\begin{ExerciseBlock}
\ExerciseHeading{}{Verständnisfragen}{differential_equations_checks}
Entscheiden Sie, welche Aussagen zutreffen, und begründen Sie Ihre Entscheidung.
\begin{enumerate}
\item Die Lösung einer Differentialgleichung ist immer eine Zahl.\TrueFalseAnswer{falsch}
\item Eine Differentialgleichung beschreibt einen Zusammenhang zwischen einer Funktion und mindestens einer ihrer Ableitungen.\TrueFalseAnswer{wahr}
\item Das zweite Newtonsche Axiom $m\ddot x(t)=F$ ist eine Differentialgleichung zweiter Ordnung für $x$.\TrueFalseAnswer{wahr}
\item Hängt die Kraft nur von der Geschwindigkeit ab, genügt für die Geschwindigkeit eine Differentialgleichung erster Ordnung.\TrueFalseAnswer{wahr}
\item Bei konstanter Kraft ist $v(t)$ eine exponentielle Funktion der Zeit.\TrueFalseAnswer{falsch}
\item Die allgemeine Lösung einer Bewegungsgleichung kann freie Integrationskonstanten enthalten und mehrere Bewegungen beschreiben.\TrueFalseAnswer{wahr}
\item Mit geeigneten Anfangsbedingungen lässt sich eine spezielle Lösung aus der allgemeinen Lösung auswählen.\TrueFalseAnswer{wahr}
\item Bei Stokes-Reibung ist der Kraftbetrag proportional zum Quadrat der Geschwindigkeit.\TrueFalseAnswer{falsch}
\item Die Gleichung $\dot v=-(b/m)v$ mit $b>0$ beschreibt einen exponentiellen Abfall des Geschwindigkeitsbetrags.\TrueFalseAnswer{wahr}
\item Bei einem exponentiellen Zerfall hängt die Halbwertszeit vom Anfangswert ab.\TrueFalseAnswer{falsch}
\end{enumerate}
\end{ExerciseBlock}
